\documentclass[11pt,a4paper]{article}
% Build: latexmk main.tex (the latexmkrc selects XeLaTeX + Biber).
% fontspec uses the bundled, static Latin Modern fonts by default.
% They support Vietnamese Unicode without system-specific Noto variable fonts.
\usepackage{fontspec}
\usepackage{geometry}
\geometry{left=2.5cm,right=2.45cm,top=2.7cm,bottom=2.6cm}
\usepackage{microtype,amsmath,amssymb,mathtools,booktabs,tabularx,array}
\usepackage{graphicx,tikz}
\usetikzlibrary{arrows.meta,positioning}
\usepackage{xcolor}
\definecolor{ink}{HTML}{16334B}
\definecolor{accent}{HTML}{116A74}
\usepackage{hyperref}
\hypersetup{colorlinks=true,linkcolor=ink,citecolor=accent,urlcolor=accent,
 pdftitle={Terrain-Aware Path Planning for Wheeled Mobile Robots}}
\usepackage[font=small,labelfont=bf]{caption}
\usepackage{enumitem}
\setlist[itemize]{leftmargin=*,topsep=2pt,itemsep=2pt}
\setlist[enumerate]{leftmargin=*,topsep=2pt,itemsep=2pt}
\usepackage{fancyhdr}
\pagestyle{fancy}
\fancyhf{}
\fancyhead[L]{\footnotesize Terrain-aware path planning for UGV}
\fancyhead[R]{\footnotesize Literature review 2026}
\fancyfoot[C]{\thepage}
\setlength{\headheight}{14pt}
\usepackage{titlesec}
\titleformat{\section}{\Large\sffamily\bfseries\color{ink}}{\thesection}{0.7em}{}
\titleformat{\subsection}{\large\sffamily\bfseries\color{ink}}{\thesubsection}{0.6em}{}
\usepackage[backend=biber,style=numeric-comp,sorting=none,maxbibnames=9]{biblatex}
\addbibresource{references.bib}
\renewcommand{\arraystretch}{1.25}
\newcolumntype{Y}{>{\raggedright\arraybackslash}X}
\newcommand{\R}{\mathbb{R}}

\begin{document}
\begin{titlepage}
\vspace*{1.7cm}
{\sffamily\color{accent}\large BÁO CÁO TỔNG QUAN NGHIÊN CỨU\par}
\vspace{1.1cm}
{\sffamily\bfseries\color{ink}\fontsize{23}{30}\selectfont
Lập kế hoạch đường đi có xét địa hình\par
cho robot di động bánh xe\par}
\vspace{0.4cm}
{\sffamily\fontsize{16}{22}\selectfont
Terrain-Aware Path Planning for Wheeled Mobile Robots\par}
\vspace{0.7cm}
{\color{accent}\rule{\textwidth}{1.3pt}}
\vspace{0.35cm}

{\large\textbf{Nền tảng thuật toán, nghiên cứu 2024--2026 và định hướng luận văn thạc sĩ}\par}
\vspace{1cm}
\begin{tabular}{@{}ll@{}}
Người thực hiện: & \textit{[Điền họ và tên]}\\
Đơn vị / lớp: & \textit{[Điền thông tin]}\\
Người hướng dẫn: & \textit{[Điền thông tin]}\\
Ngày cập nhật: & 08/10/2026\\
Phiên bản: & 1.0 --- Tổng quan ban đầu
\end{tabular}
\vfill
{\small\color{ink}Đây là báo cáo tổng quan và đề xuất phương pháp thực nghiệm.\\
Chưa có kết quả thực nghiệm do người viết thực hiện; các research gap là giả thuyết cần xác minh.}
\end{titlepage}

\begin{abstract}
\noindent
\textbf{Tóm tắt.}
Robot bánh xe tự hành (wheeled unmanned ground vehicle, UGV) hoạt động trên
địa hình phi cấu trúc cần lựa chọn đường đi không chỉ ngắn mà còn an toàn,
khả thi về chuyển động và thích ứng khi thông tin bản đồ thay đổi.
Báo cáo này giới thiệu các khái niệm terrain traversability, 2.5D elevation map,
path planning và motion planning; phân tích các nhóm thuật toán A*, D* Lite,
RRT* và Hybrid A*. Tiếp theo là tổng hợp những công trình 2024--2026 về
rủi ro địa hình, nhận thức đa cảm biến, heuristic học được và ràng buộc
động học. Dựa trên đối chiếu tài liệu, báo cáo đề xuất các câu hỏi nghiên cứu
về hiệu chuẩn rủi ro, online replanning và traversability phụ thuộc hướng,
cùng với bộ baseline và quy trình đánh giá có thể tái lập. Mục tiêu của báo cáo
là định hướng nghiên cứu thạc sĩ, không tuyên bố một thuật toán hoặc kết quả mới.

\medskip\noindent\textbf{Từ khóa:} UGV; uneven terrain; off-road planning;
terrain traversability; risk-aware path planning; D* Lite; Hybrid A*.
\end{abstract}
\tableofcontents
\clearpage

\section{Đặt vấn đề}
\subsection{Động lực nghiên cứu}
Bản đồ chiếm chỗ 2D đủ dùng cho nhiều bài toán tránh vật cản trong nhà,
nhưng chưa phản ánh được khả năng di chuyển trên địa hình dốc, gồ ghề,
có bậc cao hoặc bề mặt dễ trượt. Đường ngắn nhất có thể đi ngang sườn dốc
lớn, khiến robot mất ổn định hoặc vượt quá khả năng bám của bánh xe.
Các survey của Papadakis \cite{papadakis2013terrain} và Borges cùng cộng sự
\cite{borges2022survey} cho thấy đánh giá terrain traversability
là cầu nối giữa dữ liệu cảm biến, hình học địa hình và hệ thống lập kế hoạch.

Các nghiên cứu gần đây đã mở rộng từ cost map hình học sang lập kế hoạch
xét động học, rủi ro và tính không chắc chắn của môi trường
\cite{benrabah2024risk,wang2026winter}.
Vì vậy, chỉ cộng thêm một trọng số độ dốc vào A* không đủ để chứng minh
đóng góp học thuật mới.

\subsection{Mục tiêu và phạm vi}
Báo cáo nhằm: (i) phân biệt thuật ngữ; (ii) mô tả nền tảng
toán học và thuật toán; (iii) so sánh các công trình liên quan gần đây;
(iv) xây dựng câu hỏi và phương pháp kiểm chứng cho một đề tài thạc sĩ.
Đối tượng trọng tâm là robot bánh xe trên bản đồ độ cao 2.5D.
Bản tổng quan không giả định loại lái cụ thể; khi mô phỏng chuyển động
sẽ phải lựa chọn rõ giữa differential-drive, skid-steer và Ackermann.

\subsection{Path planning, motion planning và traversability}
\textbf{Path planning} tìm một đường hình học nối trạng thái đầu và đích.
\textbf{Motion planning} còn phải xét hướng, thời gian, vận tốc,
động học hoặc động lực học để bảo đảm đường đi có thể được thực thi.
\textbf{Traversability estimation} đánh giá địa hình phù hợp đến đâu
với một robot và trạng thái hoạt động xác định.

Ba khối này có quan hệ phụ thuộc nhưng không thay thế cho nhau.
Ví dụ một đường đi có độ dốc thấp vẫn có thể không khả thi
nếu xe không đủ bán kính quay ở một khúc cua. Ngược lại,
một vùng có bề mặt phẳng không nhất thiết có hệ số ma sát đủ lớn.

\section{Biểu diễn địa hình và mô hình toán}
\subsection{Bản đồ độ cao 2.5D}
Bản đồ độ cao có thể mô tả bởi hàm:
\begin{equation}
 z=h(x,y),\qquad h:\R^2\rightarrow\R.
 \label{eq:dem}
\end{equation}
Trong trường hợp rời rạc, mỗi ô có độ phân giải không gian \(\Delta\)
và một hoặc nhiều thống kê như độ cao trung bình, phương sai, độ dốc.
2.5D có ưu điểm gọn nhẹ nhưng không biểu diễn đúng các bề mặt
chồng lên nhau như cầu vượt và hang. Point cloud/voxel 3D phù hợp hơn
với môi trường có cấu trúc phức tạp nhưng đòi hỏi xử lý bổ sung.

\subsection{Slope, roughness, step và footprint}
Nếu bề mặt khả vi, độ dốc tổng hợp:
\begin{equation}
 S(x,y)=\arctan\!\sqrt{
  \left(\frac{\partial h}{\partial x}\right)^2+
  \left(\frac{\partial h}{\partial y}\right)^2}.
 \label{eq:slope}
\end{equation}
Theo hướng chuyển động \(\theta\), các đạo hàm có hướng:
\begin{align}
 s_{\parallel}(x,y,\theta)
    &=\nabla h(x,y)^\top(\cos\theta,\sin\theta)^\top,\\
 s_{\perp}(x,y,\theta)
    &=\nabla h(x,y)^\top(-\sin\theta,\cos\theta)^\top.
\end{align}
Chúng có thể làm đặc trưng gợi ý dốc dọc và dốc ngang,
nhưng không bằng trực tiếp pitch/roll của thân xe.
Góc thân còn phụ thuộc kích thước, hệ treo và vị trí bánh.

Một chỉ số độ gồ ghề cục bộ:
\begin{equation}
 R=\sqrt{\frac{1}{|\mathcal N|}\sum_{i\in\mathcal N}
    (z_i-\widehat z_i^{\mathrm{plane}})^2},
 \label{eq:rough}
\end{equation}
với \(\widehat z_i^{\mathrm{plane}}\) là độ cao của mặt phẳng
khớp trong vùng lân cận. \emph{Step height} có thể gần đúng bởi
độ chênh cao lớn nhất trong một cửa sổ. Các đại lượng này nhạy với
độ phân giải, nhiễu và kích thước cửa sổ.
Ngoài hình học, tính khả thi phụ thuộc đường kính bánh,
khoảng sáng gầm, tải trọng, cơ cấu lái và độ bám.

\subsection{Pipeline nhận thức đến lập kế hoạch}
\begin{figure}[htbp]
\centering
\begin{tikzpicture}[
 >=Stealth, node distance=6mm,
 every node/.style={font=\small},
 block/.style={draw=ink,rounded corners=3pt,fill=ink!5,
  text width=11.2cm,align=center,minimum height=0.84cm},
 focus/.style={draw=accent,very thick,rounded corners=3pt,fill=accent!8,
  text width=11.2cm,align=center,minimum height=0.9cm}]
\node[block] (a) {DEM / LiDAR / camera và vùng chưa quan sát};
\node[block,below=of a] (b) {Trích xuất slope, roughness, step height, uncertainty};
\node[focus,below=of b] (c) {Robot-aware traversability, giới hạn an toàn và risk map};
\node[block,below=of c] (d) {Global/local planner: A*, D* Lite, Hybrid A*, RRT*};
\node[block,below=of d] (e) {Kiểm tra khả thi / mô phỏng / cập nhật bản đồ};
\foreach \x/\y in {a/b,b/c,c/d,d/e}{\draw[->,thick,accent](\x)--(\y);}
\end{tikzpicture}
\caption{Sơ đồ khái niệm của pipeline terrain-aware navigation,
không phải một thuật toán mới.}
\label{fig:flow}
\end{figure}

\section{Các thuật toán nền tảng}
\subsection{Dijkstra và A*}
Dijkstra tìm đường có tổng chi phí nhỏ nhất trên đồ thị hữu hạn
với trọng số không âm. A* \cite{hart1968formal} sử dụng:
\begin{equation}
 f(n)=g(n)+h_{\rm heur}(n),
\end{equation}
với \(g\) là chi phí từ đầu đến nút \(n\), \(h_{\rm heur}\)
là ước lượng chi phí còn lại. Với heuristic admissible
và các điều kiện triển khai thích hợp, A* bảo đảm tính tối ưu
trên \emph{đồ thị đang giải}; không có bảo đảm an toàn vật lý
nếu bản đồ hoặc mô hình xe không đúng.

\subsection{D* Lite và Field D*}
D* Lite \cite{koenig2002dstarlite} tái sử dụng thông tin tìm kiếm
khi chi phí trên đồ thị thay đổi, phù hợp với bản đồ được cập nhật
trong quá trình di chuyển. Field D* \cite{ferguson2006field}
dùng nội suy để giảm phụ thuộc vào các hướng đi rời rạc
của lưới và cải thiện hình dạng đường đi.

\subsection{RRT* và Hybrid A*}
RRT* và PRM* thuộc họ lấy mẫu với các kết quả quan trọng
về tính tối ưu tiệm cận \cite{karaman2011optimal}.
Hybrid A* tìm kiếm trên trạng thái có hướng,
sử dụng motion primitives để phản ánh giới hạn lái.
Nhìn chung, Hybrid A* có thể đánh giá tính khả thi chuyển động
tốt hơn A* trên ô lưới 2D, đổi lại là không gian tìm kiếm lớn hơn.

\begin{table}[htbp]
\centering\small
\caption{Vai trò và giới hạn của các baseline được đề xuất.}
\begin{tabularx}{\textwidth}{@{}p{2.0cm}YYY@{}}
\toprule
Planner & Ưu điểm & Giới hạn & Nên dùng để \\
\midrule
A* & Dễ cài đặt, kết quả giải thích được & Phụ thuộc biểu diễn lưới & Benchmark path cost \\
D* Lite & Tận dụng tìm kiếm cũ khi map thay đổi & Cần cập nhật chi phí cạnh đúng & Benchmark replanning \\
RRT* & Linh hoạt trong không gian liên tục & Phụ thuộc số mẫu và steering & So sánh sampling-based \\
Hybrid A* & Xét heading và giới hạn lái & Trạng thái lớn; cần mô hình xe & Kiểm tra kinematic feasibility \\
\bottomrule
\end{tabularx}
\end{table}

\section{Tổng quan nghiên cứu liên quan}
\subsection{Các survey quan trọng}
Papadakis \cite{papadakis2013terrain} hệ thống hóa phương pháp
phân tích traversability; Borges và cộng sự \cite{borges2022survey}
phân tích các đặc trưng cảm biến, thuật toán và tập dữ liệu.
Benrabah và cộng sự \cite{benrabah2024risk} chỉ ra nhiều cách
định nghĩa rủi ro và sự cần thiết phải đánh giá rủi ro trên
toàn đường đi. Wang và cộng sự \cite{wang2024survey}
tổng quan path planning trong môi trường phi cấu trúc.
Điểm rút ra là traversability luôn phải xét trong bối cảnh
robot, dữ liệu và nhiệm vụ cụ thể.

\subsection{Các nghiên cứu trực tiếp trên UGV}
Zhang và cộng sự \cite{zhang2022wheeled} sử dụng bản đồ có xét
độ dốc, bậc và độ gồ ghề, A* cho global path và Q-learning
cho tránh chướng ngại cục bộ khi bản đồ chưa biết đủ.
Lee và cộng sự \cite{lee2023uncertainty} xem xét uncertainty-aware
navigation trên địa hình 3D. Wang và cộng sự \cite{wang2024history}
đề xuất cây tìm kiếm kết hợp đồ thị lịch sử để đi trên địa hình
chưa biết trước. Vì vậy, riêng ý tưởng đưa uncertainty
vào planner không còn mới.

Feng và cộng sự \cite{feng2025constraint} kết hợp traversability,
Patch-RRT* và tối ưu quỹ đạo có ràng buộc trong môi trường công trường.
Ji và cộng sự \cite{ji2025nnpp} nghiên cứu mô hình heuristic dựa trên học máy
để thu hẹp không gian tìm đường. Beycimen và cộng sự
\cite{beycimen2026fusion} kết hợp thông tin hình học và bề ngoài
để tạo traversability map thời gian thực. Công trình này
xuất bản trực tuyến cuối 2025 và được xếp trong tập tạp chí năm 2026.

Gai và cộng sự \cite{gai2026framework} kết hợp thông tin địa hình,
giới hạn động học và độ trơn của chuyển động trong framework
điều hướng cho xe mặt đất. Wang và cộng sự \cite{wang2026winter}
mô hình hóa các sự cố trượt, lật và lún trong điều kiện mùa đông
và tích hợp đánh giá rủi ro với Hybrid A*.

\begin{table}[p]
\centering\footnotesize
\caption{Literature matrix. Cột cuối là gợi ý câu hỏi tiếp nối,
không phải khẳng định các tác giả đã thừa nhận một limitation.}
\label{tab:lit}
\begin{tabularx}{\textwidth}{@{}p{2.25cm}p{2.7cm}p{3.05cm}Y@{}}
\toprule
Công trình & Bài toán & Phương pháp & Câu hỏi cần tìm hiểu tiếp \\
\midrule
Zhang et al. 2022 \cite{zhang2022wheeled}
& Địa hình biết một phần & A* + Q-learning & Rủi ro vùng chưa quan sát?\\
Lee et al. 2023 \cite{lee2023uncertainty}
& Điều hướng 3D bất định & Learning-based uncertainty & Hiệu chuẩn với rủi ro vật lý?\\
Wang et al. 2024 \cite{wang2024history}
& Điều hướng chưa biết map & Tree + history graph & Thay đổi địa hình ảnh hưởng latency?\\
Feng et al. 2025 \cite{feng2025constraint}
& Địa hình công trường & Patch-RRT* + tối ưu quỹ đạo & Bền vững khi sai số bản đồ?\\
Ji et al. 2025 \cite{ji2025nnpp}
& Tăng tốc tìm đường & Heuristic học được & Chi phí và giới hạn an toàn?\\
Beycimen et al. 2026 \cite{beycimen2026fusion}
& Traversability mapping & Geometry + appearance fusion & Uncertainty của map đến planner?\\
Gai et al. 2026 \cite{gai2026framework}
& Động học và terrain & Planning + tracking & Kiểm chứng dưới rủi ro bất định?\\
Wang et al. 2026 \cite{wang2026winter}
& Rủi ro mùa đông & Risk model + Hybrid A* & Chuyển miền sang mặt đất khác?\\
\bottomrule
\end{tabularx}
\end{table}

\section{Định hướng bài toán và phương pháp khởi đầu}
\subsection{Đầu vào, đầu ra và giả định}
Cho elevation grid map \(\mathcal{M}\), trạng thái đầu \(s\),
đích \(g\), đặc trưng hình học và thông số robot.
Phiên bản 1 tìm đường \(\pi=(v_0,\ldots,v_N)\) trên lưới.
Phiên bản 2 có thể tìm kiếm trên trạng thái
\((x,y,\theta)\), đồng thời kiểm tra bán kính quay và tư thế xe.
Không nên gọi phiên bản 1 là motion planning đầy đủ.

\subsection{Hàm chi phí minh họa}
Với cạnh có chiều dài thực \(d_{ij}>0\), một baseline đơn giản:
\begin{equation}
 c_{ij}=d_{ij}\left(
 1+\lambda_s \widehat S_{ij}
  +\lambda_r \widehat R_{ij}
  +\lambda_u \widehat U_{ij}\right),
 \quad \lambda_s,\lambda_r,\lambda_u\geq 0.
 \label{eq:cost}
\end{equation}
Các đại lượng có dấu mũ được chuẩn hóa trên cùng quy ước
và không âm. Thành phần 1 cho phép dùng khoảng cách Euclid
như một cận dưới của tổng chi phí trong điều kiện phù hợp.
\(\widehat U\) là độ bất định của phép đo hoặc mô hình;
nó \textbf{không đồng nghĩa} với xác suất robot thất bại.

Các điều kiện khả thi ban đầu:
\begin{equation}
 \pi\subseteq\mathcal X_{\mathrm{free}},\quad
 S_{ij}\leq S_{\max},\quad
 H^{\mathrm{step}}_{ij}\leq H_{\max}.
\end{equation}
Các ngưỡng phải liên hệ với robot và mô hình an toàn cụ thể.

\subsection{Hướng phát triển: chance-constrained risk}
Nếu có thể ước lượng xác suất sự cố đáng tin cậy,
một hướng phát triển là:
\begin{equation}
 \begin{aligned}
 \min_\pi\quad &J_{\mathrm{nominal}}(\pi)\\
 \mathrm{s.t.}\quad &
 \mathbb{P}(\text{failure trên toàn }\pi)\leq\varepsilon.
 \end{aligned}
 \label{eq:chance}
\end{equation}
Không được cộng các xác suất sự cố rồi gọi đó là xác suất
của cả đường. Bất đẳng thức hợp cho cận trên
\(\mathbb{P}(\cup_i F_i)\leq\sum_i\mathbb{P}(F_i)\)
mà không cần giả định độc lập, nhưng có thể bảo thủ.
Việc hiệu chuẩn và kiểm tra phụ thuộc giữa các sự kiện là
một phần quan trọng của nghiên cứu sâu hơn.

\subsection{Ba research questions cần kiểm chứng}
\begin{enumerate}[label=\textbf{RQ\arabic*.}]
\item \textbf{Uncertainty:} Mô hình uncertainty được hiệu chuẩn
có giúp giảm vi phạm giới hạn an toàn trên bản đồ nhiễu/thiếu dữ liệu
so với chi phí hình học đơn thuần hay không?
\item \textbf{Online replanning:} D* Lite với cập nhật cost map
có giảm độ trễ so với chạy lại A* từ đầu mà vẫn đạt cùng điều kiện
an toàn và tiêu chí cost hay không?
\item \textbf{Motion feasibility:} Bổ sung heading, turning radius
và traversability phụ thuộc hướng có giảm đường đi không thực thi được
trong mô phỏng robot hay không?
\end{enumerate}
Các câu hỏi này \emph{chưa được khẳng định là research gap mới}.
Trước khi chọn luận văn cần đối chiếu sâu với
\cite{lee2023uncertainty,wang2024history,feng2025constraint,gai2026framework,wang2026winter}.

\section{Thiết kế thực nghiệm và tiêu chí đánh giá}
\subsection{Baseline}
Thứ tự ưu tiên: (1) A* khoảng cách; (2) A* với slope/roughness
và ràng buộc cứng; (3) A* cộng uncertainty penalty;
(4) D* Lite trên cùng bản đồ chi phí;
(5) Hybrid A* khi đã xác định rõ loại lái.
Cần giữ cố định độ phân giải map, kích thước robot,
điểm đầu--đích và ngân sách thời gian tính toán.
Không so sánh trực tiếp giá trị tối ưu của hai planner
giải hai bài toán khả thi khác nhau.

\subsection{Kịch bản và chỉ số}
Sử dụng nhiều loại DEM: phẳng, đồi dốc, gồ ghề,
và vùng chưa quan sát; thay đổi mức nhiễu và số lần cập nhật bản đồ.
Tách dữ liệu xây dựng cost map khỏi dữ liệu ground truth dùng
để kiểm tra độ an toàn.

\begin{table}[htbp]
\centering\small
\caption{Chỉ số cần báo cáo cho thí nghiệm.}
\begin{tabularx}{\textwidth}{@{}p{3.2cm}Y@{}}
\toprule
Chỉ số & Ý nghĩa / điều kiện \\
\midrule
Planning success rate & Tỷ lệ trả về đường; khác với tỷ lệ robot thực thi thành công \\
Path length (m) & Độ dài trong hệ tọa độ vật lý, không chỉ số ô lưới \\
Terrain cost & Dùng cùng định nghĩa cost khi so sánh \\
Constraint violations & Vượt ngưỡng dốc, bậc, roll risk hoặc clearance \\
Runtime / expansions & Báo cáo máy chạy, số node mở rộng, trung vị và phân vị cao \\
Replanning latency & Thời gian phản ứng khi map cập nhật, với cùng ngân sách tính toán \\
Failure / slip / rollover & Chỉ kết luận khi mô phỏng/đo đạc đủ tin cậy \\
Risk calibration & Mức khớp giữa xác suất dự đoán và tần suất quan sát \\
\bottomrule
\end{tabularx}
\end{table}

\subsection{Ablation và reproducibility}
Tiến hành ablation bỏ lần lượt slope, roughness,
uncertainty và giới hạn chuyển động.
Thử nhiều random seed và nhiều vùng địa hình khác nhau;
báo cáo độ phân tán hoặc khoảng tin cậy,
không chỉ lấy một hình đường đi đẹp.
Công bố tham số xe, map resolution, seed, thư viện,
phần cứng và mã xử lý dữ liệu.

\section{Research gap và kế hoạch thạc sĩ}
\subsection{Phân biệt limitation với khoảng trống thực sự}
Một tác giả chưa triển khai một chức năng không chứng minh
rằng cộng đồng chưa giải quyết nó. Research gap cần có
ba loại bằng chứng: (i) literature search đủ sâu;
(ii) thất bại hoặc giới hạn đo được của baseline mạnh;
(iii) phương pháp mới có thể khắc phục và được kiểm chứng độc lập.
Một công thức trọng số cộng hoặc thay A* bằng Hybrid A*
không tự tạo ra tính mới đủ cho một tạp chí Q1.

\subsection{Ba hướng đề tài ứng viên}
\begin{description}[leftmargin=0pt,style=nextline]
\item[\textbf{A. Risk-calibrated path planning.}]
Tập trung vào cách chuyển uncertainty thành xác suất rủi ro
có thể kiểm định và tối ưu đường với giới hạn xác suất sự cố.
Đối chiếu trực tiếp \cite{benrabah2024risk,lee2023uncertainty,wang2026winter}.
\item[\textbf{B. Incremental replanning under uncertain updates.}]
Kết hợp cập nhật map có sai số với tìm kiếm gia tăng,
kiểm tra độ trễ và chi phí tính toán khi thông tin mới đến;
A* khởi động lại và D* Lite là baseline.
\item[\textbf{C. Direction-dependent vehicle-aware planning.}]
Đánh giá rủi ro leo/xuống/ngang dốc trên \((x,y,\theta)\),
gắn với bán kính quay và ổn định của xe;
so sánh \cite{feng2025constraint,gai2026framework}.
\end{description}

\subsection{Rủi ro của luận văn}
Nghiên cứu trên DEM tự sinh không thể chứng minh xe thật
không trượt hay lật. Semantic labels cũng không trực tiếp là
ground truth về độ ma sát. Nếu hướng tới công bố ở tạp chí
robotics có yêu cầu cao, cần dùng benchmark độc lập,
mô phỏng đã kiểm tra và thực nghiệm thực khi có điều kiện.
Kết luận về Q1 phải xét năm và hệ thống phân hạng cụ thể,
không thể bảo đảm từ tên đề tài.

\subsection{Lộ trình sáu tháng}
\begin{table}[htbp]
\centering\small
\caption{Lộ trình đề xuất để chuyển từ literature review sang thực nghiệm.}
\begin{tabularx}{\textwidth}{@{}p{1.9cm}Yp{4.1cm}@{}}
\toprule
Giai đoạn & Công việc & Bàn giao \\
\midrule
Tháng 1 & A*, D* Lite, DEM, slope, roughness & Baseline A* và báo cáo đọc bài\\
Tháng 2 & Cost map và benchmark nhiều địa hình & Kết quả có thể tái lập\\
Tháng 3 & Nhiễu, uncertainty, risk calibration & Bằng chứng cho vấn đề nghiên cứu\\
Tháng 4 & Thiết kế phương pháp và protocol & Bản mô tả thuật toán\\
Tháng 5 & Ablation, robustness, statistical tests & Bảng kết quả và phân tích\\
Tháng 6 & Replanning hoặc motion feasibility & Bản thảo luận văn/bài báo đầu tiên\\
\bottomrule
\end{tabularx}
\end{table}

\section{Kết luận}
Terrain-aware path planning cho robot bánh xe là sự kết hợp
giữa hiểu biết địa hình, thuật toán tìm đường và mô hình
khả năng di chuyển của xe. A* là baseline phù hợp để bắt đầu,
D* Lite là nền tảng nghiên cứu cập nhật bản đồ, và Hybrid A*
mở rộng bài toán sang trạng thái có hướng. Các công trình
2024--2026 đã nghiên cứu sâu về risk, traversability
và kinematics; vì vậy, nghiên cứu mới cần một vấn đề
được chứng minh bằng so sánh có kiểm soát.
Hướng ưu tiên là lập kế hoạch có xét rủi ro/uncertainty
trên elevation map và đánh giá khả năng cập nhật đường đi.
Phạm vi luận văn cuối cùng chỉ nên chốt sau khi đã tái hiện
baseline và kiểm chứng giả thuyết research gap.

\clearpage
\printbibliography[heading=bibintoc,title={Tài liệu tham khảo}]
\end{document}
